\small
\textbf{Abstract.}
A retailer makes decisions from the same demand at different grains and on
different cycles: store orders by item and day, promotions by item and week in
blocks of weeks, category reviews week by week, and supplier ranges and terms
over a year. It wants one forecast that serves all of them without losing value
between plans. This paper asks whether one forecast, built at the finest grain
and summed to every coarser grain, can serve these decisions across every kind
of line, and which kind of model does that best. It compares gradient-boosted
trees (LightGBM), pretrained foundation models (TimesFM, Chronos-Bolt,
Timestral), an attention-based model (the Temporal Fusion Transformer) and a
world model that carries each item's state forward under the price and
promotion plan, on four public data sets (M5, six Dominick's categories,
Breakfast at the Frat and Corporaci\'on Favorita), with the world model and
LightGBM tuned by the same search. Results are read by grain, by the nature of
the line, on seen and unseen data, and across retailers. On lines and stores the
models were trained on, the tuned world model and LightGBM are close, each ahead
on different data sets and grains, and their average is the most accurate
choice across grains: it is in the 90\% model confidence set at \rqMcsAverage{} of \rqMcsCells{}
data-set-and-grain pairs, for the fast lines that carry most sales and for the
slow tail. On items kept out of training, the world model loses at most \rqCostWmItemHi{} in
scaled error and forecasts within \rqWmFaItemOff\% of actual sales, while LightGBM loses up
to \rqCostLgbItemHi{} and forecasts as little as \rqLgbFaItemLoFrac{} of actual sales, which pulls the
average below the world model alone on all \rqKeptSets{} data sets tested. Settings tuned
on one retailer's data carry to another's, and training on other retailers' data
improved neither model by more than \rqPoolGainMaxTwo{}. On these data, the rule that follows is
one forecast for every item in every store, the average of the tuned world model
and tuned trees summed to every grain, with the world model alone for items the
models have not been trained on. The rule holds for the data sets studied; anyone
adopting the framework should rerun every comparison, including those that lost
here, on their own data. The study covers forecasts of 8 weeks and 28 days and
the price and promotion levers; whether the models hold over 26- and 52-week
plans is left for work with a retailer's own data.
% Numbers from results/analysis_tuned (nine data sets) and the transfer runs.

\vspace{4pt}
\textbf{Keywords:} retail forecasting; forecast grain; hierarchical evaluation;
forecast combination; promotions; new products; world models; foundation
models.
